\subsection{良序集的同构}

定理3. 设E和F是两个良序集。则以下两个命题中至少有一个成立：

\begin{enumerate}
\def\labelenumi{(\arabic{enumi})}
\item
  存在唯一的从 \(\mathbf{E}\) 到 \(\mathbf{F}\) 的一个段上的同构;
\item
  存在唯一的从 \(\mathbf{F}\) 到 \(\mathbf{E}\) 的一个段上的同构.
\end{enumerate}

设 \(\mathfrak{F}\) 为E的子集到F的映射的集合，其中每个映射定义在E的一个段上，并且是该段到F的一个段的同构。则该集合 \(\mathfrak{F}\) ，按关系“ \(v\) 是 \(u\) 的延拓”（在 \(u\) 与 \(v\) 之间）排序，是归纳的。因为若 \(\mathfrak{G}\) 是 \(\mathfrak{F}\) 的一个全序子集，则映射 \(u \in \mathfrak{G}\) 的定义域的并集 S 是 E 的若干段的并集，因此它本身是 E 的一个段。若 \(v\) 是 \(\mathfrak{G}\) 在 \(\Phi(E, F)\) 中的最小上界（第4号，例2），则 \(v(S)\) 是映射 \(u \in \mathfrak{G}\) 的值域的并集，因此是F的一个段。最后，对于S中的每一对元素 \(x\) , \(y\) （满足 \(x < y\) ），存在 \(u \in \mathfrak{G}\) ，其定义域包含 \(x\) 与 \(y\) 两者（因为 \(\mathfrak{G}\) 是全序的）；并且由于 \(v(x) = u(x) < u(y) = v(y)\) ， \(v\) 是S到 \(v(S)\) 上的一个同构，于是我们的断言得证。

现在令 \(u_0\) 为 \(\mathfrak{F}\) 的一个极大元（第4号，定理2），并令 \(S_0\) 为 \(E\) 的、作为 \(u_0\) 的定义域的那个段。如果我们证明要么 \(S_0 = E\) ，要么 \(u_0(S_0) = F\) ，定理即得证。我们用反证法，假设 \(S_0 \neq E\) 并且 \(u_0(S_0) \neq F\) 。那么将存在一个元素 \(a \in E\) 和一个元素 \(b \in F\) 使得 \(S_0 = ]\leftarrow, a[\) 并且 \(u_0(S_0) = ]\leftarrow, b[\) （第1号，命题1）。将 \(u_0\) 延拓为一个映射 \(u_1\) （段 \(]\leftarrow, a]\) 到 \(F\) 中的映射），置 \(u_1(a) = b\) ；由于 \(u_1\) 是 \(]\leftarrow, a]\) 到段 \(]\leftarrow, b]\) 上的一个同构，这与 \(u_0\) 在 \(\mathfrak{F}\) 中的极大性矛盾.

定理3中断言的唯一性是以下引理的推论：

引理4. 设E、F为两个良序集，且f、g为从E到F的两个递增映射，使得 \(f(\mathbf{E})\) 是F的一段，且g严格递增；则 \(f(x) \leqslant g(x)\) （对于所有 \(x \in E\) ）.

相反地，假设这些元素 \(y \in E\) （满足 \(f(y) > g(y)\) ）所成的集合非空；那么这个集合将有一个最小元素 \(a\) 。如果 \(x < a\) ，那么我们就有 \(f(x) \leqslant g(x) < g(a) < f(a)\) ，因为 \(g\) 是严格递增的。由于 \(f(E)\) 是 \(F\) 的一段, 存在 \(z \in E\) 使得 \(g(a) = f(z)\) ; \(f\) 是递增的，因此 \(f(z) < f(a)\) 蕴含 \(z < a\) 。因此

\[
f (z) \leqslant g (z) <   g (a) = f (z),
\]

这是荒谬的。

推论1. 良序集E到其自身一个段的唯一同构映射，是E到自身的恒等映射。

在定理3中取 \(\mathbf{F} = \mathbf{E}\) 。

推论2. 设E、F为两个良序集。若存在一个从E到F的某个段T上的同构f，以及一个从F到E的某个段S上的同构g，则必有S = E，T = F，且g与f互为逆映射。

因为 \(g \circ f\) 是从 \(\mathbf{E}\) 到段 \(g(\mathbf{T}) \subset \mathbf{S}\) （\(\mathbf{E}\) 的段）上的同构；由推论1， \(g(\mathbf{T}) = \mathbf{S} = \mathbf{E}\) ，且 \(g \circ f\) 是 \(\mathbf{E}\) 上的恒等映射。类似地， \(f \circ g\) 是 \(\mathbf{F}\) 上的恒等映射，结论得证。

推论3. 良序集E的每个子集A都同构于E的一个段。

根据定理3，只需证明不存在从E到A的形如 \(S_{a}\) 的段上的同构 g。如果有的话，g 将是 E 到 E 中的一个严格递增映射，使得 \(g(a) \in S_{a}\) ，换句话说，使得 \(g(a) < a\) ；但这个不等式与引理4矛盾（其中 f 为恒等映射）。

